% LỜI TỰA --- Quyển X
\chapter*{Lời Tựa}
\addcontentsline{toc}{chapter}{Lời Tựa}
\markboth{Lời Tựa}{Lời Tựa}

\begin{truyen}{Lời Tựa}{The Illusions of Youth}
\chuhoa{G}{iáo viên bắt đầu bộ sách này với một câu hỏi đơn giản: làm thế nào để mỗi buổi học có thể để lại thứ gì đó đáng nhớ hơn là bài thi?}
\textit{(A teacher began this series with a simple question: how can each class leave behind something more memorable than a test?)}

Ảo tưởng không phải điều xấu~--- nó là một phần của tuổi trẻ. Điều xấu là không nhận ra nó kịp.
\textit{(Illusions are not bad things~--- they are part of being young. What is bad is not recognizing them in time.)}

Quyển mười đặt tên cho những ảo tưởng phổ biến nhất của tuổi trẻ: ảo tưởng về tài năng, về tình yêu, về thành công nhanh, về việc mình đặc biệt hơn người khác. Không nhằm dập tắt ước mơ~--- mà nhằm giúp người đọc nhìn thật hơn trước khi đặt cược cả cuộc đời.
\textit{(Volume ten names the most common illusions of youth: illusions about talent, love, fast success, and being more special than others. Not to crush dreams~--- but to help readers see more clearly before betting everything.)}

Bộ sách <<Truyện Tích Cảm Hứng Song Ngữ>> gồm 28 quyển, mỗi quyển một chủ đề riêng. Quyển X này mang chủ đề: \textbf{nhận ra và điều chỉnh ảo tưởng trước khi đặt cược cả cuộc đời}.
\textit{(The series ``Inspiring Bilingual Tales'' contains 28 volumes, each with its own theme. Volume X focuses on: \textbf{nhận ra và điều chỉnh ảo tưởng trước khi đặt cược cả cuộc đời}.})
\end{truyen}

\begin{baihoc}
Nhận ra ảo tưởng của mình không phải là thất bại~--- đó là bước đầu tiên của sự trưởng thành.
\textit{(Recognizing your own illusions is not failure~--- it is the first step toward maturity.)}
\end{baihoc}

\begin{ghinhoanh}
\textbf{Short English to remember:}

Seeing clearly is the beginning of growing up.

Most illusions dissolve only under honest light.
\end{ghinhoanh}

\begin{tuhocnhanh}
1. Em đang mang ảo tưởng nào mà chưa chịu nhìn nhận? \textit{(What illusion are you holding onto without admitting it?)}

2. Điều gì trong quyển này khiến em cảm thấy hơi khó chịu? Tại sao? \textit{(What in this volume made you slightly uncomfortable? Why?)}

3. Sau khi đọc quyển này, em sẽ thay đổi cách nhìn gì về bản thân? \textit{(After reading this volume, what will you change about how you see yourself?)}
\end{tuhocnhanh}

\begin{center}
\textit{\color{truyengold}``Thức tỉnh muộn còn hơn không bao giờ thức tỉnh.''}
\end{center}

\begin{flushright}
\textbf{Tôi Nguyễn Văn Sang}\\
\textit{GV THPT Nguyễn Hữu Cảnh - TP HCM}
\end{flushright}

\ngancach
